\section{BIRD 执行轨迹与补充分析}\label{app:bird-traces}
本附录展开正文的补充执行诊断。历史记录中的 Phase I 指 BIRD 发现研究，Phase II 指独立的 BIRD 准入研究；它们均不作为 MATH 结果的因果解释。
\subsection{发现研究的程序群体与条件}\label{app:phase1-provenance}

\paragraph{程序群体与评分。}
\S\ref{sec:f1-collapse} 的发现比较在 $60$ 道题上评估由 GLM-5.3-Flash 构建器通过 TTHE 提议者生成的十二个候选，以及 \bare{} 和 \react{}。准入检查导入与冒烟执行，不验证所声称的机制。基线在发现阶段的两次运行之间翻转了 $60$ 题中的三题，因此两次运行的基线错题集略有不同。发现研究使用行集合代理评分器；在两个保留了最终 SQL 的矩阵上，它与官方 BIRD 评分器的一致率为 0.97--0.99，没有代理判对而官方判错的案例。

\paragraph{开发暴露与记录边界。}
人工控制是在查看发现轨迹后设计的，其 $151$ 个评估项中有 $24$ 个与协议开发测量重叠。因此，其余量是描述性的，不是稳定互补性的盲估计。早期会话内生成记录没有保留审核等生成预算所需的生成工具与逐尝试日志。这些限制适用于此处保留的发现证据。

\begin{figure}[H]
\centering
\includegraphics[width=0.60\linewidth]{figures/fig_phenomenon.pdf}
\caption{\textbf{源码差异未必改变观测正确性。}BIRD 发现研究包含 $14$ 个 harness、$60$ 个任务，每个执行一次。含基线的 $91$ 对中，token 序列源码相似度与正确性分歧的秩相关为 $\rho=-0.010$；$36/91$ 对具有相同正确性向量。这是观测任务上的结果比较，不是程序等价检验。}
\label{fig:phenomenon}
\end{figure}
\FloatBarrier
\subsection{发现研究的结果与人工控制}\label{app:bird-discovery}

每个 BIRD 发现候选都通过代码差异检查。然而，$21/66$ 个生成候选对具有完全相同的正确性向量（含基线为 $36/91$），$91$ 对上源码相似度与正确性分歧几乎没有秩相关（$\rho=-0.010$；图~\ref{fig:phenomenon}）。整体平均分歧为 $3.3\%$（仅生成候选为 $3.8\%$），对 \bare{} 错题的联合修复为 $7.7\%$，单次余量为 $3.3$\,pp。在这 $60$ 题上一致不能确立程序普遍等价，却说明有必要检查程序实际做了什么。


\subsection{轨迹区分机制缺失、未触发与执行失败}
\textbf{T1：描述存在，实现缺失。}十二个候选中有六个在文档字符串中声称检索或自检，却只调用一次模型、不执行任何产物，所命名的机制并未实现。\textbf{T2：机制存在，未被触发。}全部十个经过审核的 \react{} 任务中，首个 SQL 均成功执行，已有的执行—反馈—修复分支从未启动，最终 SQL 与 \bare{} 相同。执行成功不等于语义正确，这些运行并未检验失败条件下的修复。\textbf{T3：执行发生，输出损坏。}十二个候选中有一个执行多轮查询与修复，却将自然语言作为最终 SQL 返回，丢掉了 \bare{} 能答对的题。

这些观察定位不同的失败环节，而不是赋予一个共同的多样性分数。T1 要求检查实现，T2 要求任务能够触发分支，T3 要求有效的输出契约。它们未确立机制触发会引起准确率提升。四个人工编写、执行路径不同的控制在 $60$ 题上产生 $6.67$\,pp 单次余量，在 $151$ 题上为 $7.28$\,pp，说明评估装置能够记录不同修复。但这些控制是在查看轨迹后编写的，缺少匹配重复与克隆对照。附录~\ref{app:phase1-provenance} 披露较早的代理评分器及与协议开发重叠的 $24/151$ 题；附录~\ref{app:examples} 提供具体轨迹。


\FloatBarrier
\subsection{轨迹审核的完整样例}\label{app:examples}

前述 T1 计数来自对全部十二个发现候选的分类。独立的轨迹审核覆盖五个 harness：\bare{}、\react{}、单次调用候选、超时重试候选和多轮修复候选。它展示具体执行路径，不是“十二个中有六个属于 T1”的计数分母。本节将一项任务沿这些路径展开。产物逐字取自该审核（5 个 harness $\times$ 10 道题）；其日志器把存储的 SQL 摘录上限设为 120 字符，故截断以 \emph{[trunc.]} 标记。

\paragraph{任务。}BIRD 开发集问题 352（\texttt{card\_games} 数据库；题目文本逐字，拼写错误为数据集原有）：
\begin{quote}\small\ttfamily
Calculate the percentage of the cards availabe in Chinese Simplified.
\end{quote}
其 BIRD 提示：
\begin{quote}\small\ttfamily
'Chinese Simplified' is the language; percentage = Divide(Sum(id where
language = 'Chinese Simplified'), Count(id)) *100
\end{quote}
Gold SQL（在发布数据库上执行两者分别得 8.7734 与 8.7728——官方评分器将其区分）：
\begin{quote}\small\ttfamily
SELECT CAST(SUM(CASE WHEN T2.language = 'Chinese Simplified' THEN 1 ELSE 0
END) AS REAL) * 100 / COUNT(T1.id) FROM cards AS T1 INNER JOIN foreign\_data
AS T2 ON T1.uuid = T2.uuid
\end{quote}
\bare{} harness 回答：
\begin{quote}\small\ttfamily
SELECT SUM(CASE WHEN language = 'Chinese Simplified' THEN 1 ELSE 0 END) *
100.0 / COUNT(id) AS percentage FROM foreign\_data;
\end{quote}
被官方评分器判错：它直接在 \texttt{foreign\_data} 上计数，而 gold 查询在连接 \texttt{cards} 之后计数；提示没有写明该连接。

\paragraph{\react{}（T2：机制存在但未触发）。}\react{} 的控制流确实执行 SQL 并在数据库错误时重试。但其第一次调用与 \bare{} 的提示逐字节相同，而此处生成的 SQL 执行无错误，因此错误触发的重试从不发生，最终 SQL 与 \bare{} 相同。在全部 10 道被审任务上模式都成立：有 SQL 执行，10/10 最终 SQL 与 \bare{} 逐字节相同——反馈循环是真实代码，但其修复分支在这些题上未激活。需要受控故障注入来检验反馈是否改变输出。

\paragraph{单次调用候选（T1：机制被描述但未执行）。}该候选的文档字符串宣称一种“保守修复”：限制求解器最大输出 token 并要求仅答案块。其行为在结构上就是 \bare{}：一次 LLM 调用，然后提取 SQL。token 上限仅作为环境默认值存在（2{,}048 对运行时的 32{,}000）——没有任何 harness 控制流实现它，且在这些题上它从不生效，因为 SQL 答案比上限低几个数量级。在该样例题上它恰好发起一次调用、不执行 SQL，并返回与 \bare{} 相同的错误查询（忽略空白差异；直接在 \texttt{foreign\_data} 上计数而无连接）。在全部 10 道被审任务上它记录的 SQL 执行数为零，且最终 SQL 有 6 个与 \bare{} 逐字节相同；其余差异来自改写的系统提示，没有实现检索、自检或修复操作。提示改写本身是一种已实现的变换，但没有实现文档宣称的额外机制。因此最常见的 Phase-I 失败是：声明的操作没有进入实际执行的程序。

\paragraph{超时重试候选（代码级多样性，行为同一性）。}该候选是一个真正不同的程序：它把单次 coder 调用包在守护线程看门狗中，若调用停滞超过 40\,s 就发出一次对冲重复追问。机制存在于控制流中——但其触发器（停滞的 API 调用）在审核中从未出现：零次 SQL 执行，且全部 10 道被审题的最终 SQL 都等于 \bare{}。即使新增机制是真实的，代码不同而行为相同仍可存活：它只是从不触发。

\paragraph{多轮修复候选（T3：多轮路径执行失败）。}唯一一个实现了真正执行—检查—修复的被审候选把模型的\emph{自然语言说明}当作最终 SQL 返回。在一道 \bare{} 答对的题（“State the alternative languages available for card named Annul numbered 29.”）上，其记录的最终 SQL 逐字开头为：
\begin{quote}\small\ttfamily
with "only a single read-only SELECT is allowed". The previous "query" was
actually just prose text (the model wrote rea\emph{[trunc.]})
\end{quote}
修复提示导致输出契约失败：被要求修复一个“实际上只是自然语言说明”的查询时，冻结目标继续返回自然语言推理；由于找不到 SQL 块，推理文本便被当作查询提交。该题从 single-call 正确翻转为 harness 错误。


\FloatBarrier
